\section{Cote 149, scans 79--85: Kummer theory and geometric tangential splittings}
\label{sec:archive-kummer}

\paragraph{Evidence and notation.}
All seven supplied TeX transcriptions were read in full.  The original
manuscript images were not supplied, so the editorial conclusions below are
mathematical readings, not verified readings of handwritten glyphs.  Write
\(D\) for the boundary of a smooth proper relative curve \(X/S\), put
\(U=X\setminus D\), and let \(g_i:S\longrightarrow D\) be a boundary section.
This avoids the manuscript's reuse of \(T\), which we reserve for a Tate module.
Write \(E_S=\pi_1(S)\), with geometric base points suppressed.

The passage studies a proposed cusp-normalizer extension
\begin{equation}\label{eq:arc-normalizer}
 1\longrightarrow T\longrightarrow N(L_i)
   \longrightarrow E_S\longrightarrow 1,
 \qquad T=\widehat{\mathbb Z}(1),
\end{equation}
and its obstruction class \(c_i\).  The geometric interpretation proposed in
scan 81 is the Kummer first Chern class of the conormal or normal line of
\(g_i(S)\).  The author explicitly leaves that sign convention unsettled.
Scan 85 then asks for an intrinsic profinite characterization of the
\emph{geometric} splittings inside the set of all splittings.  The obstruction
to the existence of a splitting and the recognition of that subset are
different questions.

In the calculations below, \(\mathbb Z_\ell(1)\) means the inverse system
\((\mu_{\ell^r})_r\), and \(H^j_{\mathrm{cont}}\) means continuous
cohomology of that inverse system.  We do not silently replace it by
ordinary cohomology of the inverse-limit sheaf on the small etale site.
For the full Tate module, the corresponding index set is all positive
integers ordered by divisibility.

\paragraph{F149-KUM: exact limiting Kummer sequences (proved).}
Let \(S\) be connected, normal and noetherian, and suppose that \(\ell\)
is invertible on \(S\).  Set
\[
 A=\Gamma(S,\mathcal O_S)^\times,\qquad
 P=\operatorname{Pic}(S),\qquad
 B=H^2_{\mathrm{et}}(S,\mathbb G_m).
\]
For an abelian group \(M\), put
\[
 M^\wedge_\ell=\varprojlim_r M/\ell^rM,
 \qquad
 T_\ell M=\varprojlim_r M[\ell^r],
\]
where the second transition map is multiplication by \(\ell\);
for \(A\), written multiplicatively, the first quotient is
\(A/A^{\ell^r}\).  There are natural short exact sequences
\begin{align}
 0&\longrightarrow A^\wedge_\ell
   \longrightarrow \varprojlim_r H^1_{\mathrm{et}}(S,\mu_{\ell^r})
   \longrightarrow T_\ell P\longrightarrow 0,\label{eq:arc-kum-h1}\\
 0&\longrightarrow P^\wedge_\ell
   \longrightarrow \varprojlim_r H^2_{\mathrm{et}}(S,\mu_{\ell^r})
   \longrightarrow T_\ell B\longrightarrow 0.\label{eq:arc-kum-h2}
\end{align}
Moreover,
\[
 H^1_{\mathrm{cont}}(S,\mathbb Z_\ell(1))
 \simeq \varprojlim_r H^1_{\mathrm{et}}(S,\mu_{\ell^r}).
\]
If \(P[\ell^\infty]\) has bounded exponent, then \(T_\ell P=0\), the
system \(H^1_{\mathrm{et}}(S,\mu_{\ell^r})\) is Mittag--Leffler, and
\begin{align*}
 H^1_{\mathrm{cont}}(S,\mathbb Z_\ell(1))&\simeq A^\wedge_\ell,\\
 H^2_{\mathrm{cont}}(S,\mathbb Z_\ell(1))
 &\simeq \varprojlim_r H^2_{\mathrm{et}}(S,\mu_{\ell^r}).
\end{align*}
In particular the extra hypothesis holds when \(P\) is finitely generated.

\paragraph{Proof.}
The etale Kummer sequence and Hilbert 90 give
\begin{align*}
 0&\longrightarrow A/A^{\ell^r}
   \longrightarrow H^1_{\mathrm{et}}(S,\mu_{\ell^r})
   \longrightarrow P[\ell^r]\longrightarrow 0,\\
 0&\longrightarrow P/\ell^rP
   \longrightarrow H^2_{\mathrm{et}}(S,\mu_{\ell^r})
   \longrightarrow B[\ell^r]\longrightarrow 0.
\end{align*}
The transition from level \(r+1\) to level \(r\) is induced by raising
roots of unity to the \(\ell\)-th power.  In the Kummer diagram the first
copy of \(\mathbb G_m\) maps by that power and the second by the identity.
Consequently the transitions on the left are quotient maps, and the
transitions on the right are multiplication by \(\ell\).  Both left
systems are surjective, so their first derived inverse limits vanish.
Taking inverse limits proves
\eqref{eq:arc-kum-h1}--\eqref{eq:arc-kum-h2}, including the terminal zeros.

If \(\ell^bP[\ell^\infty]=0\), the transition
\(P[\ell^m]\longrightarrow P[\ell^r]\) is zero whenever \(m\geq r+b\).
The image of \(H^1(S,\mu_{\ell^m})\) at level \(r\) is then contained in
the unit term \(A/A^{\ell^r}\).  It contains that entire term because
the unit quotient transitions are surjective.  These stabilized images
prove the Mittag--Leffler assertion and also give \(T_\ell P=0\).
For a countable inverse system, the continuous-cohomology comparison is
\begin{equation}\label{eq:arc-jannsen}
 0\longrightarrow \varprojlim\nolimits^1_r
 H^{j-1}_{\mathrm{et}}(S,\mu_{\ell^r})
 \longrightarrow H^j_{\mathrm{cont}}(S,\mathbb Z_\ell(1))
 \longrightarrow \varprojlim_r
 H^j_{\mathrm{et}}(S,\mu_{\ell^r})\longrightarrow 0.
\end{equation}
The groups \(H^0(S,\mu_{\ell^r})\) are finite: a section is determined
at the generic point of the connected normal scheme and satisfies a
polynomial of degree \(\ell^r\).  Their inverse system is therefore
Mittag--Leffler.  This removes the degree-one derived-limit term without
a Picard hypothesis.  The preceding calculation removes the
degree-two term under the bounded-torsion hypothesis.  This proves all
the assertions.

The finite-level inputs are \emph{SGA 4}~\cite{SGA4}, Expose IX,
Theorie de Kummer 3.2 and Theoreme 3.3; the primary French reedition was
checked at printed pp.~14--15.  Formula \eqref{eq:arc-jannsen} is Jannsen~\cite{Jannsen1988},
\emph{Continuous etale cohomology}, \emph{Math.\ Ann.} 280 (1988),
equation (3.1), p.~216; the inverse-system Kummer diagram is
equation (3.27), p.~222.  Both locators were checked in the primary text.

\paragraph{Full completion and its precise scope.}
If \(S\) is a scheme over \(\mathbb Q\), all positive integers are
invertible.  The same proof, on the cofinal sequence \(n=r!\), gives
\begin{align*}
 0&\longrightarrow A^\wedge
 \longrightarrow H^1_{\mathrm{cont}}(S,\widehat{\mathbb Z}(1))
 \longrightarrow \varprojlim_n P[n]\longrightarrow 0,\\
 0&\longrightarrow P^\wedge
 \longrightarrow \varprojlim_n H^2_{\mathrm{et}}(S,\mu_n)
 \longrightarrow \varprojlim_n B[n]\longrightarrow 0,
\end{align*}
where \(A^\wedge=\varprojlim_n A/A^n\) and
\(P^\wedge=\varprojlim_n P/nP\).  If the torsion of \(P\) has bounded
exponent, the same stabilization argument applies to the factorial
system and identifies the middle degree-two group with continuous
cohomology.  This is the pro-\(\mathbb N\) completion.  It need not be
a compact profinite group when the quotients \(A/A^n\) are infinite.

Thus the uncertain terminal arrow in scan 82, formula (21), can be
resolved mathematically as an arrow to zero with the usual transitions.
The \(K(\pi,1)\) hypothesis in scan 80 does not by itself justify dropping
every derived inverse limit.  In the finite-generation setting the
argument above supplies the missing justification.  These observations
are not a claim to have visually read the uncertain arrow.

The requirement that the relevant integers be invertible is substantive.
For example, over \(k=\mathbb F_p(t)\), the element \(t\) acquires no
\(p\)-th root in any separable algebraic extension: such a root would
generate the nontrivial purely inseparable extension defined by
\(X^p-t\).  Hence \(x\mapsto x^p\) on \(\mathbb G_m\) is not an epimorphism
on the etale site of \(k\).  Kummer theory for noninvertible exponents
uses the flat topology; its finite flat torsors cannot simply be
reinterpreted as etale fundamental-group torsors.  In particular, the
example over \(\operatorname{Spec}\mathbb Z[1/\ell]\) below is a
\emph{pro-\(\ell\)} example, not an unqualified full-\(\widehat{\mathbb Z}\)
etale calculation on that mixed-characteristic scheme.

\paragraph{F149-SPLIT: the two torsors (proved).}
Fix a split extension \(1\to T\to E\to G\to1\), with \(T\) abelian,
and choose a continuous splitting \(s_0\).  Every other splitting is
uniquely \(s(g)=a(g)s_0(g)\), where
\[
 a(gh)=a(g)\,{}^ga(h).
\]
This is exactly the continuous cocycle equation.  Conjugating a
splitting by an element of \(T\) changes \(a\) by a coboundary.
Consequently the \emph{conjugacy classes} of splittings form a torsor
under \(H^1_{\mathrm{cont}}(G,T)\).  Raw splittings are described by
\(Z^1_{\mathrm{cont}}\), not by \(H^1_{\mathrm{cont}}\).
For the prime-invertible systems used here, the degree-one group and
etale descriptions agree: finite \(\mu_n\)-torsors are finite etale,
and their degree-zero inverse system is Mittag--Leffler as above.
No degree-two \(K(\pi,1)\) comparison is needed for this assertion.

If a line bundle \(\mathcal L\) on \(S\) is trivial, its global bases form
a torsor under \(A=\Gamma(S,\mathcal O_S)^\times\).  Indeed, two bases
\(b,b'\) have a unique ratio \(u\in A\), and every \(u\in A\) occurs.
In a local tame parameter model, replacing \(t\) by \(ut\) changes the
root tower by roots of \(u\).  Its descent cocycle is precisely the
Kummer cocycle of \(u\).  Thus the map from bases to local splitting
classes is equivariant for
\[
 A\longrightarrow H^1_{\mathrm{cont}}(S,T).
\]
The same assertion for the proposed global normalizer follows when the
relative comparison specified below has been constructed.

Higher-order parameter changes introduce no additional finite-level
class: if \(n\) is invertible in a ring \(R\), the map
\(v\mapsto v^n\) is a bijection of \(1+tR[[t]]\).
To prove this, solve successively for the coefficients of
\(v=1+\sum_{j\geq1}v_jt^j\).  At stage \(j\) the new coefficient appears
as \(nv_j\), while every other term is determined by earlier
coefficients.  The invertibility of \(n\) gives existence and uniqueness
at every stage.  A parameter change \(t\mapsto ut+O(t^2)\) therefore has
the same Kummer ambiguity as its leading unit.

\paragraph{F149-FULLINJ: injectivity of the full unit completion (proved).}
If \(S\) is an integral scheme whose function field \(F\) is finitely
generated over \(\mathbb Q\), then
\[
 \Gamma(S,\mathcal O_S)^\times
 \longrightarrow
 \varprojlim_n\Gamma(S,\mathcal O_S)^\times/
       \bigl(\Gamma(S,\mathcal O_S)^\times\bigr)^n
\]
is injective.

\paragraph{Proof.}
It suffices to prove \(\bigcap_n(F^\times)^n=\{1\}\).  Let \(k\) be the
algebraic closure of \(\mathbb Q\) in \(F\); it is a number field.
Choose a normal projective \(k\)-model of \(F\).  If \(u=v_n^n\) for
every \(n\), the order of \(u\) at every prime divisor is divisible by
every positive integer, hence is zero.  A normal noetherian domain is
the intersection, in its fraction field, of its local rings at
height-one primes.  Applying this to affine opens, both \(u\) and
\(u^{-1}\) are global regular functions on the model.  Properness then
puts \(u\) in \(k^\times\); the chosen constant field is algebraically
closed in \(F\).  Each \(v_n\), being algebraic over \(k\), also lies in
\(k^\times\).  Its finite-prime valuations show that \(u\) and every
\(v_n\) are units of the ring of integers of \(k\).  Dirichlet's unit
theorem makes this unit group finitely generated.  The intersection
of all \(n\)-multiple subgroups of a finitely generated abelian group
is zero, by its decomposition into a free group and a finite group.
Thus \(u=1\), as required.

This proves the full-completion assertion used in scan 85 under the
stated arithmetic function-field hypothesis.  It does not prove the
analogous assertion for every single prime.  For every odd \(\ell\),
\(-1=(-1)^{\ell^r}\) for all \(r\), so \(-1\) maps to zero in the
\(\ell\)-adic completion.  The full and one-prime statements must not
be interchanged.

\paragraph{F149-UNIT: a nongeometric local splitting class for every prime
(proved).}
Let \(S_\ell=\operatorname{Spec}\mathbb Z[1/\ell]\).  Then
\[
 H^1_{\mathrm{cont}}(S_\ell,\mathbb Z_\ell(1))
 \simeq
 \begin{cases}
 \mathbb Z_\ell,&\ell\ne2,\\
 \mathbb Z/2\mathbb Z\oplus\mathbb Z_2,&\ell=2.
 \end{cases}
\]
In the local model of a standard tripod cusp with a fixed parameter,
there is a local splitting class whose valuation coordinate is
\[
 \alpha_\ell=\frac{1}{1-\ell^2}\in\mathbb Z_\ell\setminus\mathbb Z.
 \tag{F149-UNIT}
\]
It is not induced by any global unit or global tangent basis, despite
\(\operatorname{Pic}(S_\ell)=0\).

\paragraph{Proof.}
The localization \(\mathbb Z[1/\ell]\) is a principal ideal domain and
has unit group \(\{\pm1\}\times\ell^{\mathbb Z}\).  Thus its Picard group
is zero, and F149-KUM identifies the displayed cohomology with the
completion of this unit group.  For odd \(\ell\), raising to
\(\ell^r\) is an automorphism of \(\{\pm1\}\); this factor disappears
in every quotient.  For \(\ell=2\) it remains a copy of
\(\mathbb Z/2\mathbb Z\), with identity transitions.  In either case
the valuation quotients are \(\mathbb Z/\ell^r\mathbb Z\), with their
usual transitions.  This proves the group formula.

The denominator \(1-\ell^2\) is a unit in \(\mathbb Z_\ell\), while
its absolute value as a reduced integer denominator is
\(\ell^2-1>1\).  Hence \(\alpha_\ell\) is \(\ell\)-adically integral
but is not an integer.  Choose integers \(a_r\) with
\(a_r\equiv\alpha_\ell\pmod{\ell^r}\).  The Kummer classes of
\(\ell^{a_r}\) are compatible, since
\(a_{r+1}-a_r\) is divisible by \(\ell^r\), and they do not depend
on the choices.  If their inverse-limit class came from a global
unit \(\pm\ell^a\), its valuation coordinate would give
\(a=\alpha_\ell\) in \(\mathbb Z_\ell\), a contradiction.

The chosen standard parameter supplies the split local tame extension.
Its splitting classes are the torsor in F149-SPLIT; changing a global
tangent basis is exactly the Kummer action of a global unit.  The
constructed class lies outside that image.  The conclusion concerns
the local splitting torsor.  It neither assumes a general relative
normalizer theorem nor proves that a richer anabelian structure cannot
recognize the geometric subset.

\paragraph{F149-TORSOR: the completion mark in scan 85.}
The group acting on \emph{bases} is the uncompleted group
\(\mathbb G_m(S)\).  Scan 82, line 4, states this correctly.  When the
Kummer map is injective, the geometric image is a sub-torsor under
that uncompleted group, whereas all splitting classes are acted on by
\(H^1(S,T)\).  F149-UNIT exhibits explicitly why replacing the former
acting group by its completion changes the statement.  Therefore the
flagged uncertain \(\wedge\) in scan 85, line 8, should mathematically
be absent if the sentence refers to the geometric bases and their
image.  Here and below an uncertainty command in the supplied source
is a source annotation, not an assertion that the original glyph was
inspected.

\paragraph{F149-CHERN-LOCAL: the finite-level line-bundle obstruction
(proved locally).}
Let \(\mathcal L=J_i/J_i^2\) be the conormal line, and let \(n\) be
invertible on \(S\).  The local root-coordinate gluing obstruction is
the Kummer boundary
\[
 \delta_n(\mathcal L)\in H^2_{\mathrm{et}}(S,\mu_n),
\]
up to inversion if the normal rather than conormal convention is used.

\paragraph{Proof.}
Choose local frames and formal parameters \(t_a\), with leading
transition units \(u_{ab}\).  The \(u_{ab}\) represent the line bundle.
On an etale refinement of the overlaps, choose
\(v_{ab}\) with \(v_{ab}^n=u_{ab}\); such refinements exist because
\(n\) is invertible.  The failure of the root-coordinate identifications
to compose on triple overlaps is
\[
 v_{ab}v_{bc}v_{ca}\in\mu_n.
\]
This is the defining boundary cocycle for the Kummer connecting map.
If necessary one uses the corresponding refinement of the Cech
hypercover; no assertion that one fixed Zariski cover computes all of
\(H^2\) is required.  The higher-order formal unit factors have unique
\(n\)-th roots by the coefficient recursion in F149-SPLIT, and do not
change this obstruction.  Dualizing \(\mathcal L\) inverts the
transition units, hence negates its cohomology class.

The local tame-cover input is \emph{SGA 1}~\cite{SGA1}, Expose XIII,
Proposition 5.2 and Corollaire 5.3, reedition pp.~315--316
(original marginal pp.~428--430).  The relative statements are
Proposition 5.5 and Corollaire 5.6.  The discrete-valuation
Abhyankar lemma is Expose X, Lemme 3.6.  The local statements and
their Kummer description were checked in the primary reedition.

\paragraph{The relative comparison and the remaining recognition gate.}
The calculation just made does not identify the complete formal-tube
extension with \(N(L_i)\) in every generality suggested by scan 79.
That identification must specify the homotopy-exactness hypotheses,
the injection and normalizer of geometric inertia, and the descent
from the local tube.  A group-cohomology comparison with
\(H^2_{\mathrm{cont}}(S,T)\) must also be supplied where it is used.
The excerpt itself marks part of that setup as requiring verification.
Theorem~\ref{thm:archive-relative} below now supplies this comparison for a smooth characteristic-zero base under the exact geometric homotopy sequence assumed in scan~79. It retains that hypothesis and does not assert the mixed-characteristic or arbitrary-fibration variants.

Once that comparison identifies the obstruction of
\eqref{eq:arc-normalizer} with the full Kummer Chern class, and once
\(\operatorname{Pic}(S)\) is finitely generated, the line bundle is
trivial if and only if that obstruction vanishes.  Indeed,
\(\bigcap_n n\operatorname{Pic}(S)=0\) by the structure theorem for
finitely generated abelian groups, and F149-KUM injects its completion
into the relevant degree-two group.  For a single prime, prime-to-\(\ell\)
Picard torsion is lost, so this full-completion conclusion cannot
simply be repeated.

A basis of the conormal line is the precise first-order datum needed
for a tangent vector.  A globally chosen formal parameter or a single
global regular function defining the divisor is a stronger datum in
general; the word \emph{uniformizer} in the excerpt should not be used
to assert such a global lifting without its hypotheses.

Over a field of characteristic zero the local splitting torsor is
already established in Stix~\cite{StixCuspidal2012}, \emph{On cuspidal sections of algebraic
fundamental groups}, author version dated 24 February 2012,
Proposition 11, p.~7.  Proposition 16, p.~8, identifies inertia
normalizers for a hyperbolic curve.  Theorem 17, p.~8, and Corollary 19,
p.~9, address injection and abundance of cusp packets under their
stated section-injectivity hypotheses.  These checked results concern
the local packet and all its section classes; they do not select the
uncompleted-unit image inside it.  Corollary 32, p.~12, recognizes
cuspidality under the non-Tate-character hypothesis, which is again
different from the requested geometric tangent subset.

\paragraph{Mathematical corrections versus unresolved readings.}
There are four evidence-based editorial entries:
\begin{enumerate}
 \item Scan 85, line 8: the acting group for global tangent bases is
 \(\mathbb G_m(S)\), as proved above; the uncertain completion mark
 should be absent in that interpretation.
 \item Scan 83, line 5: the generic points of irreducible divisors of
 \(S\) are codimension-one points of \(S\), not of the relative curve
 \(X\).  The subsequent normalization of \(S\) in its function field
 confirms this referent.  This is a mathematical correction; without
 the image it may be an author slip rather than a transcription error.
 \item Scan 82, formula (21): its terminal map is surjective and may
 end in zero by F149-KUM.  The handwritten choice has not been checked.
 \item Scan 80, line 2: a literal inference from surjectivity of a
 normalizer to the existence of a splitting is invalid.  The exact
 sequence
 \[
 0\longrightarrow\mathbb Z_\ell
 \xrightarrow{\;\ell\;}\mathbb Z_\ell
 \longrightarrow\mathbb Z/\ell\mathbb Z\longrightarrow0
 \]
 is nonsplit: a section would send a nonzero element of order \(\ell\)
 to an element of order \(\ell\) in the torsion-free group
 \(\mathbb Z_\ell\).  The following pages themselves study a splitting
 obstruction.  This proves that the unconditional implication cannot
 be right; it does not determine the replacement French wording.
\end{enumerate}

No correction of the normal/conormal sign is justified, because the
author expressly flags that convention.  Nor is the marginal extension
direction in scan 79 a mathematically compelled transcription error.
Once two genuine exact extensions have the same kernel and quotient,
a morphism inducing the identity on both is automatically an
isomorphism: its kernel is contained in the common kernel, where it
is injective, and a lift of the quotient followed by correction in the
kernel proves surjectivity.  The marginal passage discusses a setup
before all those exactness hypotheses have been established.

The Hartogs step in scan 83 needs normality of \(S\), rather than
merely normality of its open subset \(S'\).  Since the surrounding
setup already makes \(S\) smooth, the occurrence of \(S'\) does not
refute the argument and is not promoted to a required emendation.
The standard extension criterion is \emph{EGA}~\cite{EGAIV2}
\(\mathrm{IV}_2\), Theoreme 5.10.5, pp.~115--116, combined with
the normality criterion 5.8.6, p.~108, and Proposition 5.9.8 and
Definition 5.9.9, pp.~111--112.  These precise primary locators have
now been checked.  Normality gives Serre's \(S_2\) condition, so
the depth of \(\mathcal O_S\) is at least two at every point of a
closed complement of codimension at least two.  Theorem 5.10.5
then says that \(\mathcal O_S\) is closed with respect to that
complement; Proposition 5.9.8 identifies this with
\(\mathcal O_S\simeq j_*\mathcal O_{S'}\).
The equality of sections used in scan 83 follows.

\paragraph{Relation to K17.}
The archive uses Tate twist one.  Its degree-one ambiguity is a
completion of units, and its degree-two obstruction is a first Chern
class of a line bundle.  In algebraic \(K\)-theory, these are connected
with units in \(K_1\) and with the determinant or first
gamma-filtration part of \(K_0\); \(\operatorname{Pic}(S)\) should
not be inserted into \(K_0(S)\) as an additive subgroup without
qualification.

The negative-twist module
\((\mathbb Q_\ell/\mathbb Z_\ell)(-e)\), its strict local conditions,
and a duality with \(H^2(S,\mathbb Z_\ell(e+1))\) belong to a
different calculation.  A comparison must specify that change of
twist, the cohomological degree, the local conditions and the duality
map.  The Kummer sequences alone supply none of those identifications.
In particular, a bound or even a vanishing assertion for higher
\(K_{4k}(\mathbb Z)\) does not by itself select global units inside
their completion.  F149-UNIT leaves that distinction explicit even
when the Picard group is zero.

The supplied SA06 E.4 intake consists of status summaries, several
expressly marked as not read in full by its author.  It has not been
used as a proof of an all-twist cohomology table.  The relevant
war-game entry for scan 85 is S24: intrinsic recognition of the
geometric tangential subset.  The local calculations here settle the
cohomological and editorial distinctions, and Theorem~\ref{thm:archive-relative} supplies the stated relative comparison in the smooth characteristic-zero setting. The intrinsic recognition problem remains open; the two assertions are distinct.
